\chapter{Introduction}
\label{ch:into}

This chapter provides the background for the project. It introduces deep reinforcement learning, the challenges of reward scale sensitivity, our experimental configurations, and the structure of the report.

%%%%%%%%%%%%%%%%%%%%%%%%%%%%%%%%%%%%%%%%%%%%%%%%%%%%%%%%%%%%%%%%%%%%%%%%%%%%%%%%%%%
\section{Background}
\label{sec:into_back}

Reinforcement learning (RL) is a machine learning framework in which an agent learns to make sequential decisions by interacting with an environment to maximize a cumulative reward signal. Over the past decade, combining reinforcement learning with deep neural networks, known as deep reinforcement learning (DRL), has enabled performance in board games, video games, and robotics.

DRL methods can be categorized into model-free and model-based approaches. Model-free methods, such as actor-critic algorithms, learn a policy or value function directly from sampled transitions. Although widely applicable, model-free methods are often sample-inefficient, requiring millions of interactions because transitions are used only a few times for updates.

Model-based reinforcement learning (MBRL) addresses this limitation by learning a world model that predicts environment transitions and rewards. Once a world model is trained, the agent can generate synthetic experience by simulating future trajectories, augmenting its training data. This imagination-based approach can improve sample efficiency, as real transitions train a world model that generates many simulated transitions for policy optimization.

The Dreamer family of agents is widely used in imagination-based reinforcement learning. The original Dreamer agent trained actor and critic networks within imagined latent trajectories generated by a recurrent state-space model. DreamerV2 introduced discrete latent representations to improve performance on Atari benchmarks. DreamerV3 introduced innovations to enable a single set of hyperparameters to work across environments with different reward scales, observation modalities, and action spaces.

Among these innovations, the symlog transformation addresses reward scale sensitivity. The symlog function is defined as:
\begin{equation}
    \operatorname{symlog}(x) = \operatorname{sign}(x) \cdot \ln(|x| + 1),
    \label{eq:symlog_intro}
\end{equation}
with its inverse given by:
\begin{equation}
    \operatorname{symexp}(x) = \operatorname{sign}(x) \cdot \left( e^{|x|} - 1 \right).
    \label{eq:symexp_intro}
\end{equation}

\noindent By applying this transformation to the targets of the critic and reward prediction losses before computing the mean squared error, DreamerV3 compresses large magnitudes logarithmically while preserving the sign, reducing sensitivity to reward scale. DreamerV3 used this approach to achieve high performance across benchmarks without environment-specific tuning.

%%%%%%%%%%%%%%%%%%%%%%%%%%%%%%%%%%%%%%%%%%%%%%%%%%%%%%%%%%%%%%%%%%%%%%%%%%%%%%%%%%%
\section{Problem Statement}
\label{sec:intro_prob_art}

Standard mean squared error (MSE) losses are sensitive to the scale of target values. Large targets can produce large gradients, leading to unstable updates. In supervised learning, this is addressed via data normalization. In reinforcement learning, however, value estimates and predicted rewards are non-stationary signals that change during training.

This sensitivity to reward scale is a challenge for imagination-based agents. During autoregressive rollouts, prediction errors in rewards or state transitions accumulate over time steps. If these errors inflate the predicted rewards or value targets, the MSE loss amplifies the gradients, pushing value predictions further. This feedback loop can cause a value explosion, where the critic loss diverges (exceeding $10^{15}$ in our experiments), corrupting gradient signals and causing policy collapse.

The severity of this issue depends on the environment. In environments with small, bounded rewards (such as CartPole-v1, where each step yields a reward of $+1$), prediction errors are constrained, and standard MSE losses may suffice. In environments with larger, variable rewards (such as LunarLander-v3, where rewards span from approximately $-100$ to $+200$), the MSE amplification is pronounced, and training instability can emerge. Both environments are provided by the Gymnasium library~\citep{towers2024gymnasium}.

This project investigates whether the symlog transformation, originally proposed as part of the DreamerV3 architecture, can be applied in isolation to address value explosion in a simplified imagination-based agent.

%%%%%%%%%%%%%%%%%%%%%%%%%%%%%%%%%%%%%%%%%%%%%%%%%%%%%%%%%%%%%%%%%%%%%%%%%%%%%%%%%%%
\section{Aims and Objectives}
\label{sec:intro_aims_obj}

\noindent\textbf{Aim:} To evaluate whether symlog-transformed losses improve training stability and prevent value divergence in imagination-based reinforcement learning compared to standard MSE losses.

\vspace{0.3cm}
\noindent\textbf{Objectives:}
\begin{enumerate}
    \item \textbf{Implement a simplified imagination-based RL agent in PyTorch.} This agent operates in the raw state space using two-layer MLP networks with 256 hidden units for the actor, critic, and world model.
    \item \textbf{Implement the symlog transformation and integrate it into training.} Symlog is applied to the targets of the critic and reward prediction losses, computing the MSE in symlog space.
    \item \textbf{Compare three configurations:} (A) a model-free actor-critic baseline, (B) an imagination-augmented agent with standard MSE losses, and (C) an imagination-augmented agent with symlog-transformed MSE losses. These are evaluated on CartPole-v1 and LunarLander-v3 using five random seeds, yielding 30 runs.
    \item \textbf{Analyze learning curves, loss trajectories, and world model prediction errors} to assess performance and stability.
\end{enumerate}

%%%%%%%%%%%%%%%%%%%%%%%%%%%%%%%%%%%%%%%%%%%%%%%%%%%%%%%%%%%%%%%%%%%%%%%%%%%%%%%%%%%
\section{Solution Approach}
\label{sec:intro_sol}

Our approach uses a simplified Dreamer-like architecture that retains imagination-based training but operates directly in the raw observation space. This design choice isolates the effects of symlog-transformed losses by avoiding the complexity of latent-space learning.

The world model consists of a transition model predicting next-state residuals and a reward model predicting immediate rewards. Both are implemented as two-layer MLPs with 256 hidden units. The actor and critic networks use the same architecture.

Training uses a Dyna-style scheme. At each step, the agent samples transitions from a replay buffer to update the world model. It then generates imagined transitions by rolling out the world model for a fixed horizon. These real and imagined transitions are used to update the actor and critic networks.

The configurations differ only in their loss target computations:
\begin{itemize}
    \item In \textbf{Condition A} (model-free), the agent trains only on real transitions using standard MSE.
    \item In \textbf{Condition B} (imagination, standard MSE), the agent trains on real and imagined transitions using standard MSE.
    \item In \textbf{Condition C} (imagination, symlog MSE), the critic and reward prediction targets are symlog-transformed before computing the MSE loss.
\end{itemize}

\noindent All other hyperparameters are held constant. Each configuration is evaluated over five independent random seeds to account for stochastic variability.

%%%%%%%%%%%%%%%%%%%%%%%%%%%%%%%%%%%%%%%%%%%%%%%%%%%%%%%%%%%%%%%%%%%%%%%%%%%%%%%%%%%
\section{Summary}
\label{sec:intro_summary}

This chapter outlined the context and motivation for studying symlog-transformed losses in imagination-based reinforcement learning. We described the sensitivity of standard MSE losses to reward scale, defined the project aims and objectives, and outlined our simplified architecture. Chapter~\ref{ch:lit_rev} reviews the literature on these topics.